% !TeX program = xelatex
\documentclass{main}

% Replace these values with your paper metadata.
\coursetitle{《课程名称》课程论文}
\cntitle{主标题}
\papersubtitle{——副标题}
\aauthor{姓名}
\faculty{XX院校，XX系}
\stuid{学号}

\begin{document}
\makepages

\section*{引言}

这是一段课程论文正文示例，用于展示五号宋体、Times New Roman 西文、1.5 倍行距与首行缩进两字符。正文应围绕明确问题展开，交代材料范围、论证线索与核心观点。引用文献时统一使用脚注，不另设文末参考文献表。\footnote{示例作者：《中文著作示例》，示例出版社，2026年，第1页。}

\section{一级标题示例}

每节集中处理一个具体问题。篇幅约三千字时，建议正文不超过三节，不含引言与结语。论述应区分原始材料、已有研究与作者判断，避免只罗列观点。译著脚注应写明译者。\footnote{示例作者：《译著示例》，示例译者译，示例出版社，2026年，第2页。}

\section{第二个一级标题示例}

正文可以由多个自然段组成。学位论文脚注注明学位授予单位、学位类型、答辩年份与页码；期刊论文注明期刊名、年份与期号。具体著录顺序见本模板的 README 与课程格式原件。

外文专著保留原文书名，并给出出版地、出版社、年份与页码。\footnote{Example A. Author, \textit{Example Book Title}, Example City: Example Press, 2026, pp. 5--6.}

\section{第三个一级标题示例}

模板自动生成“一、”“二、”“三、”形式的一级标题。引言与结语使用不编号的 \verb|\section*|。若材料来自网页，应提供可核查的页面信息、访问日期和完整地址，并确认链接能够正常换行。\footnote{Example Organization, “Example Web Page,” [2026-07-31]. \url{https://example.com/a/long/reference/address}.}

\section*{结语}

结语应回应引言提出的问题，概括论证所得，不引入未经讨论的新材料。使用模板时，请删除全部示例文字并替换为自己的论文内容，同时保留课程名、题目、作者、院系与学号等元数据命令。

\end{document}
